% !TEX root = ../main.tex
\section{Method}
\label{sec:method}
\paragraph{Design grammar.} Every effect is measured on items that hold content fixed and vary only the target
attribute, so content-driven activation cancels in the difference. Each attribute is tested in an
\emph{explicit} (stated) and a \emph{proxy} (inferred) encoding, crossed in a $2\times2\times2$ factorial per
domain. The attribute is placed in two ways, with different roles.

\emph{Cross-marker decision design (the harm evidence).} The user turn presents one record --- a loan
application, a biography, an essay --- in each of the eight factorial cells and asks for the domain's binary
decision; the assistant turn is one of a fixed set of responses that state no attribute value (approve with
merit reasons; decline with a neutral reason; decline with a stereotype-congruent reason; decline citing the
protected attributes as categories; an evasive answer), so the identical response is valid under every cell.
Per cell the reward model's decision margin is $D = r(\text{approve}) - r(\text{decline})$, and the disparity
$D(\text{protected}) - D(\text{reference})$, averaged over the other two factors, is how much more the model
favours declining one group than another on the identical document --- what a policy optimised against it
would learn. The full $2^3$ decomposition of $D$ gives the intersectional and non-additive terms. Adapting
\citet{kumar2025prefix}, \crossinf{} is the change in decision accuracy (approve a strong record, decline a
weak one) between protected and reference cells, also in a threshold-free form (the AUC of $D$ between strong
and weak records), a reliability harm that maps onto the Act's concern with biased false positives and
negatives. The record is the unit of analysis: every contrast is taken within a record, and intervals come
from a bootstrap over records.

\emph{Direct scoring (the mechanism layer).} The rendered document is scored as the assistant's response to a
neutral assessment request, and \autoinf{} \citep{kumar2025prefix} measures the reward shift between two
byte-identical texts that differ in one marker clause. This establishes that the reward is sensitive to
protected attributes under controlled substitution, not that the model assesses applicants in a biased way,
and it yields the sharpest directions for the mechanistic analysis.

\paragraph{Mechanistic step.} The reward is linear in the pooled final-token state, $r = w\cdot h + b$, so
projecting a unit direction $u$ out of the state changes the decision disparity by exactly
$-(w\cdot u)(\Delta_\text{int}\cdot u)$, where $\Delta_\text{int}$ is the marker $\times$ decision interaction
of the states. A direction removes the disparity only if the head reads it \emph{and} it carries the
interaction. We estimate \diffmean{} directions from both placements --- the direct arm's (marker in the
response), and, cross-fitted over record folds, the prompt-placement attribute direction, the interaction
direction and an ``unfair decision'' direction (attribute-citing vs.\ neutral decline) --- project each out,
sweep the strength $\alpha$, and compare the directions' geometry against split-half reliability. A bias is
\emph{reward-scalar} low-complexity if nulling removes it at near-zero cost to legitimate accuracy. Whether a
direction fitted with the marker in the response removes the disparity with the marker in the prompt is a
transfer test in its own right (RQ5), and scoring the same cells in both placements checks how much of the
direct-scoring effect is an artefact of placing the attribute in the assistant's mouth.

\paragraph{Representation-level test.} Because null-space projection and a linear probe only interrogate the
linear subspace, we additionally apply \leace{} (optimal linear erasure) and then train a non-linear (MLP)
probe on the erased representation. A concept that is linearly erasable but non-linearly recoverable is
entangled (high-complexity) regardless of what the reward scalar shows --- the discriminator at the heart of
this work.
